\documentclass{subfiles}
\begin{document}
    We now turn our attention to a more interesting question: 
        fixed a source over some alphabet, how do we design the optimal code for it?
        The question, apparently hard, has a somewhat obvious answer:
        take in account the lengths of the codewords. Or better said, 
        their average.

    Formally speaking, we take in account the \gls{acl} of the code, 
        defined as the quantity
        \[
            \acl = \sum_{i = 1}^{n} p_{i}l_{i}
        \]
        such that is minimal. We say that a code \(\C*\) is compact (or optimal) for the source 
        \(S\) if, among all prefix codes for \(S\), \(\C*\) has the minimal ACL.

    \begin{remark*}
        Is there a lower bound for the ACL? Yes, and it given by the following result.
        \begin{lemma*}[Shannon]
            If \(\C*\) is a UD code for a source \(S\) with probabilities 
            \(p_{1}, \ldots, p_{n}\), and \(d\) is the size of the code alphabet,
            then
            \[
                \acl \ge \frac{H(S)}{\log d}
            \]
        \end{lemma*}
    \end{remark*}

    In addition to the ACL, two additional quantities are significant when designing a 
        code: the efficiency \(\eff\) and the redundancy \(\red\). More precisely,
        we want that the efficiency, defined as 
        \[
            \eff = \frac{H(S)}{\acl}
        \]
        is maximized; while the redundancy, defined as 
        \[
            \red = 1 - \eff = \frac{\acl - H(S)}{\acl}
        \]
        is minimized.

    A question still remains: how do we build optimal codes? 
        Over the years, several attempts to provide a procedure to build optimal codes 
        were done, but until 1952 none achived the goal. 
        We discuss this procedure in the next section.
        \clearpage
\end{document}
